\documentclass[namedate,webpdf,imaiai]{ima-authoring-template}

% The IMA class supplies the publication layout, fonts, margins, hyperlink
% treatment, and natbib. Keep shared manuscript content independent of the
% clean/commented wrapper; the wrapper controls only TODO visibility below.
\usepackage{booktabs}
\usepackage[nameinlink,noabbrev]{cleveref}

\providecommand{\PaperCommented}{0}
\newcommand{\draftnote}[1]{%
  \ifnum\PaperCommented=1
    \par\medskip\noindent\textbf{TODO:} #1\par\medskip
  \fi
}

\title[Mathematics support attendance and academic performance]{Attendance at a Mathematics Support Centre and Academic Performance in First-Year University Mathematics}
\author{Heriberto Espino-Montelongo*\ORCID{0009-0009-1230-2931}
\address{Universidad de las Am\'ericas Puebla}}
\author{Daniela Cort\'es-Toto
\address{Universidad de las Am\'ericas Puebla}}
\authormark{Espino-Montelongo and Cort\'es-Toto}
\corresp[*]{Corresponding author: \href{mailto:heriberto.espinomo@udlap.mx}{heriberto.espinomo@udlap.mx}}

% Required IMA front-matter fields. Replace the provisional volume, page, and
% DOI values with the journal-assigned metadata after acceptance.
\DOI{}
\copyrightyear{2026}
\vol{00}
\pubyear{2026}
\firstpage{1}

\abstract{
Evaluating mathematics support with administrative data is difficult because attendance is not randomly assigned, attendance records may cover several academic periods, and raw grades may not be directly comparable across classes. We analysed 6,627 students taking their first attempt at a first-year university mathematics course during periods for which attendance records from the Centro de Aprendizaje de Matem\'aticas (CMAT) were available. Visits were matched to the same academic period as the course, and final grades were standardised within instructor--period classes. Attendance was grouped as 0, 1, 2, 3, and 4 or more visits. Overall, 18.6\% of students attended CMAT at least once during the relevant academic period. Mean within-class standardised grades were $-0.061$ standard deviations among students with no recorded visits and $0.219$, $0.247$, $0.278$, and $0.355$ among students with 1, 2, 3, and 4 or more visits, respectively. Games--Howell post hoc comparisons showed that each group of attendees had a higher mean standardised grade than students with no recorded visits, whereas none of the six comparisons among attendance-frequency groups was significant at the 5\% level. Linear regression models with class fixed effects and adjustment for degree programme gave the same overall pattern: estimated differences relative to students with no visits were $0.305$, $0.359$, $0.373$, and $0.461$ standard deviations for 1, 2, 3, and 4 or more visits. The findings show a consistent association between CMAT attendance and higher academic performance within class, but much less evidence that students who attended more frequently performed progressively better. Because attendance was not randomly assigned and prior mathematical attainment was incompletely observed, the results should be interpreted as associations rather than causal effects.
}

\keywords{mathematics learning support; mathematics support centre; attendance; academic performance; first-year mathematics; higher education}

\begin{document}
\maketitle

\section{Introduction}

Mathematics learning support is now an established feature of higher education in many countries, particularly as institutions respond to differences in students' prior mathematical preparation and to the quantitative demands of a widening range of degree programmes \citep{lawson2020,mullen2024}. A substantial body of research has examined the impact and evaluation of mathematics support, but observational evaluation remains difficult because students are not randomly assigned to attend. In many settings participation is voluntary, while in others attendance may also be influenced by institutional incentives. Students who attend can therefore differ from those who do not in prior attainment, motivation, perceived difficulty, study effort, time availability, and willingness to use academic support. Self-selection is consequently a central concern in studies based on attendance records \citep{mullen2024}.

Studies of mathematics support centres have commonly compared the grades or pass rates of students who attend with those of students who do not \citep{macanbhaird2009,jacob2019}. \citet{jacob2019}, in particular, found that students who attended once should be considered separately from students who never attended and reported more favourable outcomes at higher levels of attendance. A new observational study therefore needs to do more than reproduce a general difference between attendees and non-attendees. It should also consider whether attendance is measured for the relevant course period, whether grades are comparable across classes, and whether different frequencies of attendance are associated with clearly different levels of performance.

The broader literature on academic help-seeking provides a useful framework for interpreting attendance. Help-seeking is generally understood as a self-regulated and motivated learning strategy, and its relationship with achievement depends on both the form of help sought and student characteristics \citep{fong2023}. Attendance at CMAT can therefore be viewed as an observable form of formal academic help-seeking, but the number of visits does not directly measure the amount or quality of support received. Students may return for different reasons, including continuing difficulty, persistence, confidence in the service, scheduling opportunities, or recommendations from others.

To make the attendance--performance comparison more interpretable, this study uses three features of the research design. First, CMAT visits are matched to the same academic period as the student's MU course attempt rather than accumulated across different periods. Second, final grades are standardised within instructor--period classes so that students are compared with others whose work was evaluated in the same local grading context. Third, students are grouped by 0, 1, 2, 3, and 4 or more visits rather than assuming that each additional visit has the same relationship with performance. These choices do not remove self-selection or other unmeasured differences between students, but they reduce two avoidable problems: assigning visits from other academic periods to the course result being analysed and pooling raw grades across classes with different grade distributions.

The study addresses two questions: (1) how is CMAT attendance during the same academic period as MU associated with the within-class standardised grade on a student's first MU attempt; and (2) among students who attend CMAT, do standardised grades differ clearly according to the number of visits? The analysis is observational throughout and does not attempt to estimate a causal effect of mathematics support.

\section{Institutional context}
